\noindent\textbf{(a)} \[100x''+40x'+80x=F_0\sin(4\pi t),\quad x(0)=x'(0)=0.\]

\noindent\textbf{(b)} Sea \(\omega=4\pi\), \(\nu=\sqrt{19}/5\), \(R=4/5-\omega^2\), \(Q=R^2+(2\omega/5)^2\),
\[A=\frac{R}{100Q},\quad B=-\frac{2\omega/5}{100Q}.\]
Entonces
\[x(t)=F_0\left[A\sin\omega t+B\cos\omega t+e^{-t/5}\left(-B\cos\nu t+\frac{-A\omega-B/5}{\nu}\sin\nu t\right)\right].\]

\noindent\textbf{(c)} La curva \(x/F_0\) es una oscilación rápida superpuesta a un transitorio de frecuencia \(\nu/(2\pi)\); termina oscilando con amplitud \(1/(100\sqrt Q)\).

\noindent\textbf{(d)} No existe tal frecuencia real con los datos dados: \(c^2-4km=-30400<0\). La condición no describe un experimento realizable.
